% Как один сэмпл превращается в данные для каждой из четырёх архитектур
\begin{tikzpicture}[
  font=\small,
  node distance=4mm and 8mm,
  src/.style={
    rectangle, rounded corners=2pt,
    draw=violet!60!black, fill=violet!10,
    minimum width=64mm, minimum height=12mm,
    inner sep=4pt, align=center
  },
  view/.style={
    rectangle, rounded corners=2pt,
    draw=teal!60!black, fill=teal!8,
    minimum width=42mm, minimum height=24mm,
    inner sep=4pt, align=center
  },
  ttl/.style={font=\bfseries\small},
  arr/.style={-{Latex[length=2mm]}, thick, draw=black!60}
]

\node[src] (s) {Сэмпл (мастер): \texttt{node\_features} (T,N,24) + \texttt{coords} + \texttt{edges} + \texttt{times}};

% PINN
\node[view, below left=12mm and 32mm of s] (pinn)
  {\ttl PINN\\[2pt]
   \scriptsize случайные точки\\$\{(x_i, t_i)\}$:\\
   collocation, BC, IC,\\
   data-loss};

% FNO
\node[view, below left=12mm and -2mm of s] (fno)
  {\ttl FNO\\[2pt]
   \scriptsize интерполяция\\на регулярную\\сетку\\$(N_x, N_y, N_z, N_t)$};

% MGN
\node[view, below right=12mm and -2mm of s] (mgn)
  {\ttl MeshGraphNet\\[2pt]
   \scriptsize узлы + рёбра,\\окно $t \to t{+}1$,\\автoрегрессионно};

% Transformer
\node[view, below right=12mm and 32mm of s] (tr)
  {\ttl Transformer\\[2pt]
   \scriptsize input-токены\\$(p_i, a_i, \mathrm{IC}_i)$,\\query-токены\\$(x_q, t_q)$};

\draw[arr] (s.south) -- ($(pinn.north)+(0,2mm)$);
\draw[arr] (s.south) -- ($(fno.north)+(0,2mm)$);
\draw[arr] (s.south) -- ($(mgn.north)+(0,2mm)$);
\draw[arr] (s.south) -- ($(tr.north)+(0,2mm)$);

\end{tikzpicture}
